% Full inherited normalization and kernel proofs, with the new trace interface.
% Retained normalization proof, with three explicit references redirected to the archive.
\section{Physical kernels, first-exit traces and analytic dependencies}
\label{app:v83-kernels}
This appendix supplies the measure-theoretic interfaces used by the new
arguments in one convention. It does not deduce a surface estimate from a
volume multiplier theorem. The geometric and spectral estimates whose
proofs are retained in the technical supplement are listed separately.

\subsection{Normalization and common kernels}
On a regular word use endpoint angles $q=(\alpha_0,\alpha_m)$ and
oriented internal arclengths. Endpoint first variation gives
$F_{\alpha_0}=-Rp_0$ and $F_{\alpha_m}=Rp_m$. Consequently
\[
 \omega=|\partial_{\alpha_m}p_0|=|F_{\alpha_0\alpha_m}|/R,
 \qquad d\bar\nu=\frac{\omega}{4\pi c_*}\,dq.
\]
In the optical Schur complement the internal Hessian is $D_cAD_c$,
where $A$ has diagonal $t_{i-1}+t_i+2/(Rc_i)$ and off-diagonal
$-t_i$. The two angle-to-arclength factors are $R^2$; division by
$R$ in the first variation gives $\omega=Rc_0c_m/y_m$. For one
flight $y_1=\ell_0$, and for two flights
$y_2=\ell_0+\ell_1+2\ell_0\ell_1/(Rc_1)$. Changes of tangent
orientation do not change these absolute identities.

Away from $g=0$, the coordinates $(F,\vartheta)$ have Jacobian
$\chi=e^\perp{}^{\mathsf T}Ke^\perp$ in absolute value, where
$e=g/|g|$. Thus the angular expression for the original fiber is
\begin{equation}\label{eq:v83-appendix-fiber}
 d\mu_t=\frac{\omega(t,\vartheta)}{4\pi c_*\chi(t,\vartheta)}
          \,d\vartheta.
\end{equation}
The sign of an oriented endpoint chart is irrelevant; all source
Jacobians in this appendix are positive. The original source already
contains its initial and terminal tests and is normalized only by $c_*$. 

\begin{proposition}[Common graph covariance]
\label{prop:v83-common-kernel}
At fixed $m,R,n,k$, the regular exact event has jointly measurable
coarea kernels for which, for almost every $(t,h)$ and every nonnegative
Borel graph weight $f$,
\begin{equation}\label{eq:v83-common-covariance}
 \int_{S_h(t)}f(x,h,\Phi_hx)J_h(x)\,d\mu_t(x)
 =\int_{\Phi_hS_h(t)}f(\Phi_{-h}y,h,y)\,d\mu_{t+h}(y).
\end{equation}
The same versions can be used for countably many source restrictions,
positive truncations and their increasing limits. Here $S_h(t)$ is the
set of source points whose forward lifetime is greater than $h$.
\end{proposition}
\begin{proof}
At fixed count the regular physical words are disjoint open pieces
of the collision cylinder. Their true boundaries and the removed
critical points have zero source area. Take a countable cover by
precompact flow boxes with lifted angular coordinate. Assign each
point to its first containing box to obtain a Borel disjointization.
This assignment is used only for integration: its seams do not stop
the physical flow.

On the open survivor graph the map $(x,h)\mapsto(\Phi_hx,h)$ is
injective, smooth and has physical Jacobian $J_h$. Its inverse on
its image is the backward flow. Since $F(\Phi_hx)=F(x)+h$, ordinary
change of variables followed by disintegration over $(F,h)$ gives
\eqref{eq:v83-common-covariance} on a countable determining algebra
of Borel sets. Remove the union of the resulting null sets. Uniqueness
of the kernel measures and monotone approximation extend equality to
all nonnegative Borel weights outside this one joint null set.
Countable disjoint summation gives the complete physical formula.

In endpoint coordinates the field is translation in $F$ at fixed
$\vartheta$, so its endpoint Jacobian is $\chi(x)/\chi(\Phi_hx)$.
Multiplication by the physical density ratio gives
$J_h=(\omega(\Phi_hx)/\omega(x))\chi(x)/\chi(\Phi_hx)$, exactly
the normalization in \eqref{eq:v83-common-covariance}. Nonnegative
graph selectors accommodate first-exit restrictions depending on
$(x,h)$ without new incompatible versions. Monotone convergence at
fixed count permits the actual fractional insertions. No null set
independent of the uncountable radius parameter is required: the
ordered norm is $\sup_R\esssup_t$.
\end{proof}

\subsection{True outgoing faces and exhaustion}
On a nongrazing limiting word all incidence limits are positive, so
the stationary reflection equations extend smoothly to the boundary.
This continuation is not asserted to remain the physical first-hit
trajectory beyond that boundary. On a transverse outgoing arc write
$F=b(\vartheta)$. Contraction of the area form
$\omega(4\pi c_*\chi)^{-1}dF\,d\vartheta$ with $\partial_F$
gives the positive flux in \eqref{supp-v85archive:eq:v83-surface-flux}. A first-exit
point is always counted together with its incident one-sided physical
component. Lower-dimensional and tangent strata have zero transverse
flux. Section decisions, competing obstacles and simultaneous failures
are not suppressed by this description. An exit with a vanishing
internal incidence belongs to the grazing class even when another
condition fails there simultaneously; the remaining class is nongrazing.

\begin{proposition}[Initial mass and actual first-exit flux]
\label{prop:v83-first-exit-kernel}
For a nongrazing outgoing face, let $s=F(y)-t$ and $x=\Phi_{-s}y$.
On $0<s\le h$, $s<\tau^-(y)$, the contribution to initial weighted
exit mass is
\begin{equation}\label{eq:v83-inverse-jacobian-exit}
                 \frac{W^\varepsilon(x)}{J_s(x)}\,d\zeta(y).
\end{equation}
The boundary source \eqref{supp-v85archive:eq:v83-untruncated-flux} has this same
first-exit selection with weight one on the boundary, not with the
factor in \eqref{eq:v83-inverse-jacobian-exit}. Both are defined on
common almost-everywhere starting-roof versions.
\end{proposition}
\begin{proof}
In one angular flow box put
$a(u,\vartheta)=\omega(u,\vartheta)/(4\pi c_*\chi(u,\vartheta))$.
The initial fiber mass is $W(x)a(t,\vartheta)d\vartheta$, whereas
the outgoing flux is $a(b(\vartheta),\vartheta)d\vartheta$.
The ratio of the latter density to the former unweighted density is
$J_s(x)$. This proves \eqref{eq:v83-inverse-jacobian-exit}. The
backward-lifetime condition says precisely that the segment between
the starting roof and this boundary has remained in that one-sided
event; a later re-entry does not satisfy that condition.

Use precompact regular flow boxes, with rational rectangles in
$(F,\vartheta)$, and increasing finite unions whose interiors cover
the regular event. On overlaps within a physical component their
coordinates agree up to a lifted angular constant and can be refined
into finitely many pieces. Roof-constant faces are transverse;
angle-constant faces have zero flux. The outside source plus the
first-exit source for each such domain is the integral of the
indicator that the whole forward segment is not contained in the
domain. These indicators decrease to the indicator of a true exit.
Every compact surviving segment eventually lies in an exhaustion
domain. The usual endpoint convention at lifetime exactly $h$ does
not change the joint almost-everywhere formula.

For actual nongrazing faces, exhaust the regular one-sided arcs and
their backward segments by such compact pieces. Disjoint assignment
by the incident component and first-exit graph removes chart
multiplicity. Apply the local formula and nonnegative convergence
at fixed count. These operations also justify the variation payments
in Theorem~\ref{supp-v85archive:thm:v83-two-sided-flux}. No surface limit is commuted
with the collision-normalized supremum.
\end{proof}

\subsection{The indispensable analytic chain}
Input~\ref{input:v83-local} records the complete estimates used by this
article, rather than requiring a reconstruction from historical status
files. Its common-space operator dependencies are the collision norm
definitions (Supplement equations~\eqref{supp-eq:v35-weak-norm}--
\eqref{supp-eq:v35-unstable-norm}), the full-record fixed-band
Doeblin--Fortet and spectral estimates (Supplement
Theorem~\ref{supp-thm:v40-polynomial-LY}), and the finite moving-peak
local inversion. The proof uses bounded transverse multipliers, not a
bounded strong operator of multiplication by $c^{-1}$ or $|g|^{-1}$.

For specificity, the joint endpoint/incidence estimate has leading
factor $b^{129/128}s^{17/16}$ and finite-count cap
$\min\{b\rho^{m/2},s^2\}$. The endpoint test is the particular
vector $P_b\mathbf1$ of strong norm $O(b^{7/8})$; that estimate is
not a small operator norm on arbitrary vectors. Each continued
rank-one peak is normalized using a nonzero smooth test against a
physical peripheral vector, never by division by $\Pi\mathbf1$.
The homogeneous trace estimates pass through bounded finite Cesaro
approximants in the common strong space. These conventions and the
two chronological characteristic pairings are proved in the retained
module 186, not assumed for arbitrary $L^\infty$ representatives.

Positive incidence layer cake integrates the cap first and yields
a critical-volume remainder $b^{1/2}\rho^{m/4}$. A second layer cake
against $b^{-1-\eta}$ converges at finite count only for
$\eta<1/2$, giving the actual $v^{\eta,b}$ estimate stated in
Input~\ref{input:v83-local}. The leading exponent is $129/128-\eta$.
The first inverse-gradient remainder instead depends on its lower
cutoff and is not used to remove that cutoff here. For the section
layer only the all-depth marked-strip theorem is needed; for the
incidence and clearance layers the original physical-guard local
bound is used. All of these sources keep the forward occupation
coefficient. Neither signed interface cancellation nor deletion of
an arithmetic peak enters the new estimates.

Thus the new potential and surface theorems depend on explicitly
stated volume estimates, but the surface source is bounded only after
its actual first-exit graph and optical variation have been treated.
The retained spectral and billiard estimates remain available for
specialist verification. This appendix and native source compilation
are not independent human or formal proof certification.


\subsection{The new trace interface}
The new proof uses the exact segment density and the inequality
$|\partial_u A|\le CAI_m/z$, including both the optical and angular
factors. Applying the fundamental theorem of calculus to $A^{1/p}$
requires only local absolute continuity and the actual integral of
$AI_m^p$; it is not an anisotropic multiplier assertion. The reference
interval is chosen on the longer side of each actual segment, with a
Borel tie rule. Fubini then sums its energy with multiplicity one.
The resulting finite-count bound permits the two initial thresholds
to shrink with roof width. It does not extend through the unestimated
critical/zero-residence source, which is retained in the final charge.

\subsection{The enlarged moment and weighted-fiber interface}
The norms and physical multiplier realization used in
\eqref{eq:v88-multiplier}--\eqref{eq:v88-spectral} are precisely
Supplement equations~\eqref{supp-eq:v35-weak-norm}--\eqref{supp-eq:v35-unstable-norm},
Supplement Proposition~\ref{supp-prop:v49-physical-multiplier}, and
Supplement equations~\eqref{supp-v80archive:eq:v80-physical-modulus}
and \eqref{supp-v80archive:eq:v80-spectral-input}. The fixed-band
refinement used here is quantitative in the chosen moment exponent;
it uses the physical real-power length bound, not traces of an arbitrary
bounded measurable representative. Continued projections are handled by
finite physical powers and a remainder, even when $\Pi\one=0$.

The resulting $I_m^p$ moment has no inverse-gradient insertion. Its
constants may depend on the one fixed band. This is sufficient for the
finite-count supremum in \eqref{eq:v88-expanded}; an additional
band-limit claim would require additional uniformity and is not made.
The original guard modulus is Supplement
Corollary~\ref{supp-v78archive:cor:v78-guard}. The weighted-fiber
measures use one common segment correspondence and the full variation
norm. They do not compare separately selected labels or neglect exited
mass. The earlier trace interface remains in the retained source with
its original, narrower moment range.
